
\subsubsection{6.1 Key Findings}

\textbf{Finding 1: Gradient subspace methods require sufficient accumulation density.}

Accumulating one gradient per epoch produces a subspace whose rank equals the number of epochs, not the requested SVD rank. In this experiment, 10 epochs yielded a rank-10 subspace in a 25.6-million-parameter space.

\textbf{Finding 2: The PGO hypothesis remains unverified, not refuted.}

The null result stems from measurement limitation rather than absence of the hypothesized effect. The question of whether gradient subspaces can capture simplicity bias remains open pending valid measurement.

\textbf{Finding 3: Multi-batch accumulation is necessary.}

With approximately 38 batches per epoch (4,795 samples / 128 batch size) and 10 accumulation epochs, multi-batch accumulation would yield approximately 380 gradient samples---sufficient for meaningful rank-50 subspaces.

\subsubsection{6.2 Limitations}

\textbf{The core hypothesis is unverified.} PGO cannot be claimed to work or fail because the measurement apparatus prevented valid testing.

\textbf{Only one dataset tested.} Other spurious correlation benchmarks may exhibit different gradient dynamics.

\textbf{No baseline comparison performed.} Comparison against JTT, DFR, or Group DRO awaits valid PGO implementation.

\textbf{Single random seed.} Experiments used seed 42 for all runs. While sufficient to demonstrate measurement apparatus limitation, reproducibility across seeds should be verified in future work.

\subsubsection{6.3 Broader Impact}

This work identifies a methodological pitfall in gradient-based machine learning interventions. Researchers proposing gradient subspace methods should validate accumulation density before claiming directional findings.

